\chapter{검증 근거 목록}

본문의 수치는 다음 검증 근거에서 얻었다. 각 근거는 입력 조건, 판정 기준, 결과와 실행 경과를 함께 보관하며, 코드와 테스트벤치는 같은 번호 체계로 정리했다. 판정 기준을 시험 전에 정하는 것이 원칙이며, E30 이후의 시험은 기준과 코드를 실행 전에 기록했다. E26–E29와 TB21–TB23은 기준과 결과가 같은 기록에 있어 순서를 기록으로 확인할 수 없다(4장). 회로 모델은 TB01–TB33의 33개 번호이고, 변형 TB04B를 포함하면 34개 테스트벤치다.

{\footnotesize
\begin{longtable}{L{\dimexpr 0.11\linewidth-2\tabcolsep\relax}L{\dimexpr 0.25\linewidth-2\tabcolsep\relax}L{\dimexpr 0.64\linewidth-2\tabcolsep\relax}}
\caption{검증 근거 목록(E1–E47, TB01–TB33)}\\
\toprule
\tabhead \textbf{번호} & \textbf{근거} & \textbf{내용} \\
\midrule
\endfirsthead
\toprule
\tabhead \textbf{번호} & \textbf{근거} & \textbf{내용} \\
\midrule
\endhead
\bottomrule
\endlastfoot
E1 & Stage A 스위칭 시험 & 8위상 RSC, 고정 타이밍과 영교차 추종의 부품 공차 표본, 손실·모듈 셰딩·프리차지 \\
E2 & Stage B 평균모델 정합성 & 에너지 보존, 평형점 잔차, 적분간격 수렴, 이중 구현 대조 \\
E3 & 25 kW 정상상태 전력수지 & 640–920 V 평형점의 경로전류, 직렬전압, 링크전력 \\
E4 & 요구사항 분석 & 부하 입력범위와 규격 대조, 보장 최소 배터리 전압 \\
E5–E7 & 제어기 검증·지연 여유·추가 강건성 & Tier-1·Tier-2 552건, 표본화 선형화, 복합 공차 128개 코너, 부하 계약 최악 시퀀스 \\
E8 & 부하 승인 시험 & 무작위 시나리오 200개, 메시지 지연·손실, 관리 방식 3종 비교 \\
E9 & 제어기 설계 대안 비교 & 대역폭·극 배치 변경과 전압항 포화 \\
E10 & Stage C 스위칭 시험 & 3 kW CLLC의 전압이득, 입력 변동, ZVS 여유 \\
E11 & 채택 제어기 독립 구현 & MATLAB 재구현과 1,542건 대조 \\
E12 & Stage B 스위칭 모델 & PWM·데드타임·출력용량·DAB SPS, DAB 전류 제한, 데드타임 보상, 센서 사양, 게이트 시간 분해능 \\
E13 & 기준 구조 예비 비교 & ③·③′·④의 정적 범위, Tier 시험, 감쇠기 설계 \\
E14 & Stage B 보호·기동 & 정지·기동 전이, 링크 프리차지, 센서 고장과 타당성 감시 \\
E15 & 기준 구조 손실·정격 & ①–④의 손실(프로파일 가중), 반도체 수, 스위치 VA, H5 판정, 감도 7종 \\
E16 & 충전 중 보조버스 & 정적 범위, 동시 운전, 모드 전환, 전체 순서, 스위칭 모델 \\
E17 & 보호 상태기계 결합 & 감시 회귀, R2 측정 기반 상태기계 117건, 정책 비교와 감도 \\
E18 & 팩 등급별 가치 지도 & 180–240s × B\textsubscript{420}·B\textsubscript{470}의 성립성·손실·정격, 500 V 충전기 호환 \\
E19 & 창 조정 모듈 & DAB 정격 규칙, 공차 포함 Tier, 모드 전환, 고장 강등 \\
E20 & DAB 고장 모드 변조 & SPS·EPS 정확 해석과 지령표, 강등 방식 비교 40건 \\
E21 & 고장 모드 변조의 회로 여유 & 데드타임·출력용량·역방향 도통 스위칭 정상상태, TB10 재현, 첨두 표, 강등 방식 P3s·P3t(전이 여유) \\
E22 & 조정형 Stage C 5 kW & 후보 손실 비교, 직류변압기 정상상태, 레그 출력용량 모델, 역방향 경부하 이득, 제어 시험 8종 \\
E23 & R2 통합 회로 & 기본형 버스 + 조정형 Stage C 평균모델, Middlebrook 임피던스, 강등 버스 궤적의 48 V, Stage A 경부하 영전류 스위칭 \\
E24 & Stage A 경부하 위상 셰딩 & 문턱 설계, 데드타임 전류 소멸 처리, 정상 운전점 26점, 전환 과도 10사례 \\
E25 & 산업 방식 대비 구조 비교 & 모터·인버터 승압 충전 손실(96조합), 800 V급 보조부품 구조, 손실 에너지 예시, Stage A 손실 감도 \\
E26 & 충전 중 Stage A 정지 협조 & 모듈 과부하(스위칭 모델), V\textsubscript{A} 노드 사건 32건(클램프 유무, 충전기 반응, 통신 지연) \\
E27 & 영교차 게이트 방향 반전 & 무장하지 못한 반주기 처리, TB20 사례 D 재현, 방향 반전 격자 48, 회귀 \\
E28 & Stage A 4모듈 관리기 & 모듈 추가·줄임·사전 기동, 공차 표본 3개, 정상점 36·과도 7, 재시작 C\textsubscript{fly} 오차 \\
E29 & 보호 상태기계–SALS 연동 & 부하 등급 계약 응답, 메시지 지연·손실, 연동 강등 74, 48 V 회생, 이중 고장 \\
E30 & 모듈 고장 처리 & 주행·충전 중 고장 8사례 × 공칭·공차 표본 \\
E31 & Stage A 출력용량 손실 & Q\textsubscript{oss}·V 손실, 소자 4종 설계 격자 90안, 정확해 대조, 선택 규칙 \\
E32 & A25 손실·문턱 & 25 kHz 직접 모의 손실 격자, 셰딩·관리기 문턱, 보조·충전 손실 \\
E33 & A25 동특성 & 정상점·과도·느린 충전·모듈 고장 \\
E34 & 부품 정격 1차 & 9개 부품의 데이터시트 대조 \\
E35 & A25R 설계와 판정 & 설계 단계 후보, C\textsubscript{A} 2.4 mF, 요구 안·밖 과도, 모듈 고장, 정상점 \\
E36 & A25R 영향 재확인 & 손실·기존 방식·가치 지도 재계산, 충전 중 보조버스·R2 버스·충전 중 정지·이중 고장·충전기 급정지 \\
E37 & 부품 정격 2차 & C\textsubscript{fly}·클램프·이중 고장 방전·소프트 투입 소자·접촉기·C\textsubscript{A}·L\textsubscript{r} \\
E38 & 외부 검토 대응·설계 변경 & 클램프 재선정, 고전압 팩 설정, Stage C 모델 교체, 충전 손실 직접 계산, Stage C 일시 정지, 400 V 부하 정지 \\
E39 & 클램프 스위치 재판정 & 최대 열저항, 가지 저항 섞임 \\
E40 & 규칙 L 재판정 & 클램프 재선정, 최대 연속 켜짐 시간, 문턱 오차와 오동작, P1·P6·P7·P8·C\textsubscript{A} \\
E41 & ④ 대 ⑤·⑥ 공정 비교 & 보조 강압기 이상적 셰딩, 민감도 10종, 충전 경로(전용 승압기 16–50 kHz) \\
E42 & 가치 지도 공정 기준 재계산 & ①의 보조 강압기 이상적 셰딩, 팩 등급 × 부품 등급 칸 판정, 손실 계수 민감도 7종 \\
E43 & 제출 전 외부 검토 대응 & 측정 사양 S1, 출력용량 항 스위칭 배율 민감도, 클램프 자기 가열 가정, 일시 정지 동안 48 V 배터리 \\
E44 & 제출 전 재검토 대응 & 비교 소자 E\textsubscript{oss} 적용 전압 정정의 연쇄 재계산, 클램프 문턱 구현 요구, C\textsubscript{fly} 후보값 동특성, 판정 범위 재분석 \\
E45 & 제출 전 통합 수정안 대응 & H2 두 전류 기준 재집계, 가치 지도 판정 함수 정정, 충전전력 두 정의, 180s 최저 버스, 클램프 켬·끔 지연 0.5 µs 반례, P5 공칭·모서리 평균손실, 제동 펄스열, 실행 방식과 설정 전달·엄격 비교(docs/50) \\
E46 & 제출 전 두 번째 통합 검토 대응 & Tier-2 복합 공차 +3 µs 스위칭 진단, 충전 시작 보조부하 0–25 kW 첨두, 216s 직결 충전 단자와 그림 2(d), 가치 지도 전압 경로·누설 L 정정, Stage C 52 V CV 경계·열손실, TB15·TB25·E31 재집계, 보호 에너지 포트별 값·P6·DAB 1차 첨두, 검증기 오류 주입·재개 키(docs/51) \\
E47 & 최종 제출 판정 대응 & 검증기 오류 주입과 S2 재판정, 충전 과도시험 분모·부하 스텝과 TB28 기록 재집계, 입력 해시 줄바꿈 정규화, 그림 2(a) 축 이름(docs/52) \\
TB01–TB33 & Simulink 회로 시험 33종(TB04B 포함) & 시험별 회로 구성·근거 번호·판정은 표 \ref{t:tb} \\
\end{longtable}
}
\srcnote{주: 도구는 Python이 기본이다. Numba 가속 Python: E12, E16, E21, E24, E26–E28, E30, E32, E33, E35, E36, E38, E40. MATLAB 병행: E1, E2, E10. MATLAB 형식 코드의 GNU Octave 실행: E11. ngspice 병행: E22. Simulink: TB01–TB33. E45–E47의 검증기·해시·실행 방식 세부는 저장소 docs/50–52에 있다.}
